\documentclass[a4paper,11pt]{article} % Don't change this
\usepackage{fullpage} % Don't change this

\usepackage{titling} % Don't change this
\pretitle{\vspace{-6em}\begin{center}\larger[2]\bf} % Don't change this
\postauthor{\end{center}} % Don't change this
\preauthor{\vspace{-2em}\begin{center}\smaller[1]} % Don't change this
\postauthor{\end{center}\vspace{-5em}} % Don't change this

\usepackage[shortlabels]{enumitem} % Don't change this 
\setlist{itemsep=0em,topsep=0.3em} % Don't change this

%%% Some useful packages.
\usepackage{xcolor}
\usepackage{float}
\usepackage{amsmath,amssymb,mathtools}
\usepackage{graphicx}
\usepackage{subcaption}
\usepackage[hidelinks]{hyperref}
\usepackage{relsize}
\usepackage{cleveref}
\crefname{figure}{Fig.}{Figs.}

\usepackage[backend=bibtex,natbib=true,style=authoryear-icomp,maxnames=3,giveninits=true]{biblatex}
\citetrackerfalse
\bibliography{ref} % put bib entries in ref.bib

\usepackage[linesnumbered]{algorithm2e}
\usepackage{cleveref}
\crefname{equation}{Equation.}{Equations.}
\crefname{figure}{Figure}{Figures.}

\usepackage{lipsum}
\usepackage{wrapfig}
\usepackage{tikz}

%%% Add additional packages or macros below
\usepackage{array}
\usepackage{booktabs}
\usepackage{microtype}

% Draft aid: red bold TODO markers. Delete this macro's colour (or the macro
% calls) before submission.
\newcommand{\TODO}[1]{\textcolor{red}{\textbf{[TODO: #1]}}}

\newcommand{\Kh}{\widehat{K}}
\newcommand{\Kt}{\widetilde{K}}
\newcommand{\Fnorm}[1]{\lVert #1 \rVert_F}
\newcommand{\tgrok}{t_{\mathrm{grok}}}
\newcommand{\Astar}{A_{*}}

\begin{document}

\title{Scale or Shape? Decomposing NTK Movement During Grokking}
\author{Group G2 \\
Zain Al-Saffi (s4800836) $\cdot$
Jasper Chong (s4883420) $\cdot$
Aaron D'Mello (s4890257) \\
Zac Kienzle (s4803349) $\cdot$ Samuel Rieger (s4904959)}

\date{}
\maketitle

% =====================================================================
% SCAFFOLD NOTE (delete before submission)
% Every bullet below is lifted from final_proposal.tex / the G2 context
% doc and is here as a writing prompt, not as final prose. Replace each
% itemize with paragraphs. Limits: 12 pages body (excl. refs/appendices),
% 5 pages appendices. Do not change font or page size.
% =====================================================================

\begin{abstract}
	\TODO{write last, $\approx$150 words, in prose}
	Skeleton:
	\begin{itemize}
		\item Grokking is usually explained as a lazy-to-rich transition, evidenced
			by relative kernel movement $\Fnorm{K_t-K_0}/\Fnorm{K_0}$.
		\item That statistic is not identifiable: it mixes a change in kernel
			\emph{scale} with a change in kernel \emph{shape} (\cref{eq:decomp}).
		\item We split it into $S_t$ (scale), $R_t$ (shape) and $A_t$ (centred
			alignment to the labels) and use them as timing variables on modular
			addition.
		\item Head result: the fit $\log\tgrok = c + a\log N + b\log(\eta\lambda)$,
			which separates a shape clock ($a>0$, $b\approx0$) from a weight-decay
			clock ($a\approx0$, $b\approx-1$). \TODO{our $(a,b)$ with CIs}
		\item Secondary: o (or not) of the alignment threshold $\Astar$;
			a worked cell where kernel movement is large but $A_t$ is flat.
	\end{itemize}
\end{abstract}

\section{Introduction}

% Target: ~1.5 pages. Order: phenomenon -> NTK/lazy-rich -> the conflation
% -> research question -> hypothesis -> contributions.

\paragraph{The phenomenon.}
\begin{itemize}
	\item Grokking: training loss falls early, test loss stays at chance for
		orders of magnitude longer, then falls abruptly \citep{power2022grokking}.
	\item $\tgrok$ = number of training steps from the training loss first
		crossing $\tau_{\mathrm{train}}$ to the test loss first crossing
		$\tau_{\mathrm{test}}$.
	\item Task: modular addition. Input concatenates one-hot codes of
		$a,b\in\{0,\dots,p-1\}$ (dimension $2p$); target is the one-hot code of
		$(a+b)\bmod p$; $p^2$ pairs in total.
	\item Why this task: small enough for full-batch GD (clean kernel analysis);
		the generalising solution is known (Fourier ``clock''
		\citep{nanda2023progress}; closed form for quadratic activation
		\citep{gromov2023grokking}).
\end{itemize}

\paragraph{NTK, lazy and rich.}
\begin{itemize}
	\item $\Theta_t(x,x')=\langle\nabla_\theta f(x;\theta_t),\nabla_\theta
		f(x';\theta_t)\rangle$; $K_t$ is the empirical version on a fixed probe set
		of $n$ inputs.
	\item Lazy: $\Theta_t$ stays at $\Theta_0$, network behaves as its
		linearisation \citep{jacot2018ntk}. Rich: $\Theta_t$ moves
		\citep{chizat2019lazy}.
	\item One account of grokking is a lazy-to-rich transition: a linearised
		model memorises, but generalising needs the kernel to move
		\citep{kumar2024grokking,mohamadi2024grok}.
\end{itemize}

\paragraph{The conflation (the core observation).}
\begin{itemize}
	\item ``The kernel moves'' bundles two claims: a change in \emph{scale}
		(every eigenvalue rescaled, eigenbasis fixed) and a change in \emph{shape}
		(the eigenbasis itself turns).
	\item Let $\Kh=K/\Fnorm{K}$, $S_t=\log(\Fnorm{K_t}/\Fnorm{K_0})$ and
		$R_t=1-\langle\Kh_t,\Kh_0\rangle_F$. The standard statistic decomposes as
		\begin{equation}
			\Fnorm{K_t-K_0}^2\big/\Fnorm{K_0}^2 = e^{2S_t}+1-2e^{S_t}(1-R_t).
			\label{eq:decomp}
		\end{equation}
	\item Two limits make the problem concrete: at $R_t=0$ it reads
		$(1-e^{S_t})^2$ --- a kernel whose eigenvectors never move can still report
		arbitrarily large ``movement'' by shrinking; at $S_t=0$ it reads $2R_t$.
	\item Consequence: any exponent fitted to $D=\Fnorm{K_t-K_0}/\Fnorm{K_0}$ is
		unidentified. (Derivation in \cref{sec:appendix}, not asserted in the body.)
	\item Why this is not pedantry: the theorem usually cited to link weight decay
		to kernel movement describes \emph{pure scaling}. A $k$-homogeneous network
		(scaling all parameters by $c$ scales the output by $c^k$; a bias-free ReLU
		MLP has $k$ = depth) at infinite width, MSE, $L_2$ strength $\lambda$, has
		$\Theta_t=e^{-2(k-1)\lambda t}\Theta_0$
		\citep[Thms.~1--2]{lewkowycz2020l2}, i.e. $S_t=-2(k-1)\lambda t$ and
		$R_t\equiv0$: no features learned. Yet
		\citet[App.~13, arXiv v3]{kumar2024grokking} invoke it to argue weight decay
		ends lazy training and $\tgrok\propto1/\lambda$ --- an argument they call
		speculative and do not test on the kernel.
\end{itemize}

\paragraph{Research question and hypothesis.}
\begin{itemize}
	\item \textbf{RQ.} What sets $\tgrok$? Candidates: kernel shape, kernel or
		weight scale, the weight-decay rate. Three 2026 accounts each back a
		different one (\cref{sec:related}); none measures shape change. We measure
		all three on one grid and ask which $\tgrok$ tracks.
	\item \textbf{H (shape clock).} With $Y$ the probe set's one-hot label matrix
		and $H=I-\tfrac1n\mathbf{1}\mathbf{1}^\top$, let $A_t$ be the cosine between
		$HK_tH$ and $HYY^\top H$ --- centred alignment \citep{cortes2012alignment},
		invariant to rescaling of $K_t$ by construction. H says: generalisation
		begins when $A_t$ crosses a threshold $\Astar$; $\Astar$ is shared across
		$\alpha$, $\eta\lambda$, $N$ and $p$; and $\tgrok$ is independent of $S_t$
		given $A_t$.
	\item \textbf{Discriminating prediction.} From the exact evolution
		$\dot\Theta_t=-2(k-1)\lambda\,\Theta_t+T_t+T_t^\top$
		\citep[Eq.~S5, arXiv v1]{lewkowycz2020l2}: the weight-decay term only
		rescales $\Theta_t$, so $A_t$ moves only through $T_t$. $T_t$ vanishes at
		infinite width and carries no explicit $\lambda$, so shape change is a
		finite-width effect and wider networks reshape more slowly. Fit
		$\log\tgrok=c+a\log N+b\log(\eta\lambda)$: H predicts $a>0$, $b\approx0$
		($a=1$ if the $O(1/N)$ correction scaling holds); a weight-decay clock
		predicts $a\approx0$, $b\approx-1$; two nonzero exponents mean two clocks
		and the slower one sets the time. We estimate $(a,b)$, we do not assume them.
\end{itemize}

\paragraph{Contributions.} \TODO{state honestly; overclaiming costs more under
\emph{correctness} than it gains under \emph{novelty}}
\begin{itemize}
	\item Already published elsewhere: the NTK--grokking link, the laziness knob
		$\alpha$, eNTK spectral tracking, the norm delay law, norm mediation.
	\item Ours: (i) decomposition~\eqref{eq:decomp} used as a \emph{timing}
		variable; (ii) the $(a,b)$ test separating width-driven reshaping from the
		$\lambda$-clock; (iii) a shape-based account of the norm law under MSE.
	\item Practical stake: if timing is set by kernel shape, then a quantity
		measurable during the apparently-flat phase predicts capability onset before
		any loss-based metric --- an early-warning signal for delayed generalisation.
\end{itemize}

\section{Related Work}
\label{sec:related}

% Target: ~1.5 pages. Group thematically, not paper-by-paper.

\paragraph{Grokking phenomenology.}
\begin{itemize}
	\item \citet{power2022grokking}: the original delayed-generalisation
		observation on small algorithmic datasets.
	\item \citet{nanda2023progress}: what the network actually learns (Fourier
		multiplication), and the progress-measure framing --- note $A_t$ is a
		progress measure in this sense, computable online and not flat during
		memorisation.
	\item \citet{liu2023omnigrok}: grokking beyond algorithmic data; weight-norm
		framing. \TODO{decide whether to keep --- used for context only}
\end{itemize}

\paragraph{Lazy vs.\ rich training.}
\begin{itemize}
	\item \citet{jacot2018ntk}: NTK, infinite-width limit, $\Theta_t\equiv\Theta_0$.
	\item \citet{chizat2019lazy}: laziness is controlled by output scale, not
		width alone; the $\alpha$ knob with $\eta=\eta_0/\alpha^2$.
	\item \citet{kumar2024grokking}: grokking \emph{is} the lazy-to-rich
		transition; conditions (misaligned $\Theta_0$ top eigenvectors, dataset big
		enough to generalise but not so big that train tracks test, lazy start).
		Their evidence for ``the kernel changed'' is exactly the unidentified
		statistic $D$.
	\item \citet{atanasov2022silent}: silent alignment --- eigenstructure can move
		while the kernel is still small, scale grows afterwards. So the temporal
		ordering of $S$ and $A$ is an empirical question, not definitional.
\end{itemize}

\begin{itemize}
	\item \citet[\S5]{kumar2024grokking} \emph{do} measure kernel--target
		alignment, $y^\top K_0y/(\Fnorm{K}\lVert y\rVert^2)$, citing
		\citet{kornblith2019cka} --- but at initialisation only, as a static screen
		for whether $\Theta_0$ is badly enough matched to the task for grokking to
		be possible (one of their three preconditions).
	\item It is never plotted against training step. The only time-resolved
		alignment object in the paper is the misalignment \emph{error} term in the
		loss decomposition of their analytic polynomial model, not a measurement on
		the modular-addition network.
	\item Hence no threshold crossing exists in their design and $\Astar$ is not a
		question they can pose --- it is open, not settled. Our $A_t$ is the same
		family of statistic evaluated \emph{along the trajectory} and used as a
		timing variable.
	\item Two definitional gaps when our numbers sit beside theirs:
		(i) no centring matrix --- with mean-centred labels the numerators agree,
		since $\langle HKH,yy^\top\rangle_F=(Hy)^\top K(Hy)$, but their denominator
		is $\Fnorm{K}$ rather than $\Fnorm{HKH}$, so the two differ by
		$\Fnorm{K}/\Fnorm{HKH}$, and one-hot targets (column means $1/p$) are not
		centred anyway; (ii) they evaluate $K_0$, we evaluate $K_t$.
	\item So we report the uncentred variant throughout for like-for-like Tier~0
		comparison, and keep centred $A_t$ \citep{cortes2012alignment} primary
		because it is invariant to the constant component that otherwise dominates.
	\item \TODO{verify the printed form of their CKA against the PDF --- the name
		(CKA, after \citeauthor{kornblith2019cka}, which is centred) and the printed
		formula (which is not) appear to disagree; do not rest the comparison on a
		misreading}
\end{itemize}

\paragraph{Theory constraining what the kernel can do.}
\begin{itemize}
	\item \citet{lewkowycz2020l2}: Thms.~1--2 (eigenvalues decay, eigenvectors
		static) and the exact Eq.~S5/S6 (S8/S9 in v2) --- the technical spine.
		Assumptions: MSE, $k\ge2$, gradient flow, infinite width, Dyer--Gur-Ari
		conjecture.
	\item \citet{mohamadi2024grok}: no permutation-equivariant model in the kernel
		regime achieves small population error on modular addition without a
		constant fraction of the $p^2$ pairs --- so on this task shape change is
		\emph{necessary}; scale motion provably cannot deliver generalisation.
	\item \citet{gromov2023grokking}: closed-form grokked weights for the
		two-layer quadratic-activation network. Needed for Tier~4 only.
	\item \citet{tian2025emergence}: scaling laws of feature emergence.
		\TODO{one line, or cut}
\end{itemize}

\paragraph{The 2026 dispute: three clocks.}
\begin{itemize}
	\item \citet{truong2026norm}: weight norm \emph{causally} sets the timescale.
		Clamping $\lVert W\rVert$ at $\rho\lVert W\rVert_c$ gives
		$\tgrok\propto e^{a\rho}$ with $a\approx7.5$ (MLP) vs $\approx15$
		(unnormalised attention), $R^2>0.99$; $\lVert W\rVert_c\propto p^{0.38}$
		($\lVert W\rVert_c=\{42,48,55,61,67\}$ at $p=\{29,41,59,79,97\}$);
		LayerNorm removes the dependence. Loss is cross-entropy.
	\item \citet{truong2026logit}: under CE the norm effect runs through
		\emph{logit scale}, not the norm ($R^2=0.97$); the channel is absent under
		MSE, where about half the effect is unexplained.
	\item \citet{kim2026clock}: weight decay is the clock,
		$\tgrok\propto(1-\beta)/(\eta\lambda)$ with heavy-ball momentum $\beta$.
	\item \textbf{The gap.} Every observable in the dispute is scale-like or
		optimiser-level, so none can tell a reshaped kernel from a decaying one.
		H supplies the missing variable $A_t$. Our regime (MSE + vanilla GD) is the
		one where the logit channel is absent, Thm.~2 applies verbatim, and coupled
		$L_2$ coincides with decoupled weight decay, so Kim's coupled/decoupled
		distinction collapses.
	\item \citet{cortes2012alignment}: centred kernel alignment, i.e. our $A_t$.

\end{itemize}

\section{Methods}
\label{sec:methods}

% Target: ~3 pages. This is where mathematical maturity shows: derive the
% identity, do not assert it; state estimator choices and the pre-registered
% t_grok definition precisely enough to reproduce.

\subsection{Task, architecture and baseline}
\begin{itemize}
	\item Baseline follows \citet[App.~8.3, arXiv v3]{kumar2024grokking} so that
		Tier~0 is checkable against a published figure.
	\item One hidden layer, width $N=100$; input $2p$ concatenated one-hots,
		output $p$; MSE loss (accuracy is threshold-gameable on this task);
		full-batch vanilla GD at $\eta=100$; $p=23$; $90\%$ of the $p^2$ pairs for
		training; no weight decay in the baseline (swept separately).
	
	\item Table~\TODO{ref}: full config in one table (reproducibility is a graded
		criterion; synthetic data, one script, fixed seeds).
\end{itemize}

\paragraph{Parameterisation.}
\begin{itemize}
	\item Both Tier~0 and Tier~2 use the NTK parameterisation
		$f(x;\theta)=\tfrac{1}{\sqrt N}\sum_{i=1}^{N}a_i\,
		\phi(\tfrac{1}{\sqrt D}w_i^\top x)$, $a_i,w_{ij}\sim\mathcal N(0,1)$,
		$D=2p$, no biases $\Rightarrow$ $2$-homogeneous, so $k=2$ in
		\citet{lewkowycz2020l2}.
	\item $\alpha$ sits on top of this via $\tilde f=\alpha[f(x;\theta)-f(x;\theta_0)]$,
		$\eta=\eta_0/\alpha^2$: output scale and width stay separate controls.
	\item Why it is forced: under $1/\sqrt N$, $T_t$ is the only shape channel and
		the leading finite-width correction is $O(1/N)$ (Dyer--Gur-Ari), which is
		what makes $a=1$ a sharp prediction rather than a fitted slope.
	\item Under mean-field / $\mu$P ($1/N$ prefactor) feature learning is $O(1)$ at
		every width by construction: $N$ does not control laziness,
		$\hat a\approx0$ follows trivially, and the estimate says nothing about a
		shape clock. Reporting $\hat a$ without fixing the parameterisation is
		uninterpretable.
\end{itemize}

\paragraph{Consequence for Tier~0 (flag, do not bury).}
\begin{itemize}
	\item \citet[App.~8.3]{kumar2024grokking} give width, loss, optimiser, $\eta$,
		$p$ and train fraction, but \emph{not} the parameterisation, activation or
		initialisation scale; no code accompanies the paper.
	\item Their analytic model \citep[Eq.~3]{kumar2024grokking} is mean-field,
		$f=\tfrac{\alpha}{N}\sum_i\phi(w_i\cdot x)$ --- so the modular-addition
		network may not share ours.
	\item We therefore reproduce their $\alpha$ sweep \emph{under our
		parameterisation} and gate on the ordering and approximate magnitude of
		$\tgrok$ across $\alpha$, not on step counts, which the unstated
		initialisation scale alone would break.
	\item \TODO{report whether Tier~0 reproduces under NTK parameterisation. If it
		reproduces only under the mean-field prefactor, say so and state which
		parameterisation each tier uses --- a Tier~0 holding in one and a Tier~2
		requiring the other is a discontinuity to report, not to smooth over}
\end{itemize}

\subsection{The laziness and weight-decay knobs}
\begin{itemize}
	\item Laziness: centred, rescaled predictor
		$\tilde f=\alpha[f(x;\theta)-f(x;\theta_0)]$ with $\eta=\eta_0/\alpha^2$
		\citep{chizat2019lazy,kumar2024grokking}. Both parts are mandatory: centring
		removes the initial function, the $\eta_0/\alpha^2$ rescaling keeps early
		dynamics comparable across $\alpha$. Large $\alpha$ $\Rightarrow$ lazy.
	\item $\alpha$, initialisation scale and label rescaling are one knob, so we
		sweep only $\alpha$.
	\item Weight decay is always reported as the dimensionless $\eta\lambda$ ---
		the combination the theory's timescale $t_*\propto(\eta\lambda)^{-1}$ is
		expressed in. Discrete GD gives $w\leftarrow(1-\eta\lambda)w-\eta\nabla L$,
		so stability needs $\eta\lambda<2$ ($\lambda<0.02$ at $\eta=100$).
\end{itemize}

\subsection{Kernel metrics}
\begin{itemize}
	\item Fixed probe set of $n=256$ input pairs, drawn once and frozen across
		all runs and seeds. It holds test pairs, up to half of the probe, and
		training pairs for the rest. At $p=23$ with a $90\%$ training fraction
		only 53 test pairs exist, so the probe holds all 53 of them and 203
		training pairs. Metrics are reported on the training part and on the test
		part separately, as well as pooled.
	\item eNTK via \texttt{torch.func} (\texttt{jacrev} + \texttt{vmap}),
		$256\times256$ per checkpoint, seconds to compute. The sum-over-logits
		kernel adds the kernels of the $p$ outputs,
		\[
			K_t(x,x')=\sum_{c=1}^{p}\big\langle\nabla_\theta f_c(x;\theta_t),
			\nabla_\theta f_c(x';\theta_t)\big\rangle .
		\]
		This is the trace of the $p\times p$ matrix-valued
		NTK of \citet{jacot2018ntk}. The kernel of the summed output $\sum_c f_c$
		would also add the cross terms with $c\neq c'$. The two agree in the
		infinite-width limit at initialisation \citep[Theorem~1, arXiv
		v4]{jacot2018ntk}. At finite width they differ. The first-logit kernel is a
		robustness check.
	\item All of $S,R,A$ computed on the centred kernel $\Kt=HKH$; uncentred $A$
		also reported for comparability with \citet{kumar2024grokking}. They
		evaluate it on the test set \citep[Section~5, arXiv v3]{kumar2024grokking},
		so the comparison uses the test part of the probe.
	\item Definitions (\TODO{as a small table}): $S_t=\log(\Fnorm{\Kt_t}/\Fnorm{\Kt_0})$;
		$R_t=1-\langle\widehat{\Kt}_t,\widehat{\Kt}_0\rangle_F$;
		$A_t=\langle\Kt_t,\widetilde G\rangle/(\Fnorm{\Kt_t}\Fnorm{\widetilde G})$
		with $\widetilde G=HYY^\top H$.
	\item Also logged each checkpoint: train/test MSE, $\lVert w-w_0\rVert/\lVert
		w_0\rVert$, $\lVert W\rVert$, and $y^\top K^{-1}y$ (equal to
		$\lVert w-w_0\rVert^2$ near the kernel limit, so its divergence predicts
		when linearisation must break).
	\item Checkpoint on a \emph{log-spaced} step grid: grokking spans 3--5 orders
		of magnitude in $t$, so linear checkpointing spends the budget at the wrong
		end.
	\item Sanity check before trusting $A_t$: $A_0$ stable across seeds.
\end{itemize}

\subsection{Defining $\tgrok$}
\begin{itemize}
	\item $\tgrok$ = step gap between the training loss crossing
		$\tau_{\mathrm{train}}$ and the test loss crossing $\tau_{\mathrm{test}}$.
	\item Thresholds pre-registered (fixed before any sweep ran; \TODO{state
		values and where they were frozen --- repo commit hash / OSF}) with a
		sensitivity analysis over at least three values each.
	\item A run whose test loss is still above $\tau_{\mathrm{test}}$ at the step
		budget has \emph{not grokked}. It is kept, and $\tgrok$ is recorded only as
		greater than the run length: right censoring. Dropping such runs would select
		exactly the cells that confirm any timing hypothesis.
	\item Estimation by censored regression (Tobit, or Cox on $\log\tgrok$), not
		least squares on the survivors.
\end{itemize}

\subsection{Theory: where shape change can come from}
\begin{itemize}
	\item Exact evolution $\dot\Theta_t=-2(k-1)\lambda\,\Theta_t+T_t+T_t^\top$
		\citep[Eq.~S5]{lewkowycz2020l2}; $T_t$ is the third-derivative term of
		Eq.~S6. First term is pure scale, so all shape change is carried by $T_t$,
		which vanishes with width and carries no explicit $\lambda$.
	\item Under H, $\Astar$ is a fixed target, so $\tgrok$ is a distance over a
		rate: it should grow with $N$ and stay flat in $\eta\lambda$.
	\item Caveat to state plainly: the $e^{-2(k-1)\lambda t}$ solution assumes
		gradient flow, while we use $\eta=100$ discrete GD. Eq.~S5 itself is
		discrete-agnostic; the closed-form solution is not.
\end{itemize}

\subsection{Scope: four reportable tiers (five with the theory tier)}
% Each tier is independently reportable; the project cannot end with nothing.
\begin{itemize}
	\item \textbf{0} (hard gate, week~1): reproduce the published $\alpha$ sweep
		at the stated hyperparameters. Everything downstream is uninterpretable
		without it.
	\item \textbf{1}: measure $S_t,R_t,A_t,\tgrok$ over the $\alpha\times\eta\lambda$
		grid; test whether $\Astar$ is invariant; exhibit a cell where kernel
		movement is large but $A_t$ is flat. Sufficient for a complete report alone.
	\item \textbf{2} (head claim): the $(a,b)$ fit with bootstrap intervals, plus
		the residuals where a power law fails.
	\item \textbf{3} (stretch): replicate the clamped-norm law under MSE; test
		whether $\Astar$ is invariant to $\rho$.
	\item \textbf{4} (theory, droppable): bound $\lVert T_t\rVert$ for the
		two-layer quadratic-activation network (making $k=3$), whose grokked
		solution is closed-form \citep{gromov2023grokking}, giving analytic
		$A_\infty$ and $\tgrok\sim A_\infty/\dot A$; Eq.~S5 is exact, so Tiers~1--2
		stand regardless.
\end{itemize}

\section{Experiments}
\label{sec:experiments}

% Target: ~4 pages. One figure per claim, seed error bars on everything.
% Brief warns against galleries of training curves and per-algorithm tables.

\subsection{Design and budget}
\begin{itemize}
	\item Arms (\TODO{single table}): Tier~0 repro, $\alpha\in\{0.1,0.3,1,3,10,30\}$,
		6 cells $\times$ 5 seeds; Tier~1 grid, $\alpha\times\eta\lambda$ with
		$\eta\lambda\in\{0,10^{-3},3\!\times\!10^{-3},10^{-2},3\!\times\!10^{-2},10^{-1}\}$,
		36 cells $\times$ 5; Tier~2 factorial, $N\in\{50,100,400,1600\}\times
		\eta\lambda\in\{0,10^{-3},10^{-2},10^{-1}\}$, 16 cells $\times$ 5; Tier~3
		clamp, $\rho\in\{1,1.25,1.5,1.75,2\}$, 5 cells $\times$ 10; generality,
		$p\in\{23,29,31\}\times$ train fraction $\in\{50,70,90\}\%$; CE arm at the
		baseline cell.
	\item $\approx$380 runs, $\approx$20 GPU-hours. \TODO{report actual hardware
		and wall-clock}
	\item Ten seeds on the clamp arm because the MSE effect is about half-size
		\citep{truong2026logit}; CIs, not point estimates.
	\item CE arm exists to link our constants to \citet{truong2026norm}, whose
		results are under cross-entropy.
\end{itemize}

\paragraph{Estimating $b$, and the $\eta\lambda=0$ cell.}
\begin{itemize}
	\item $\log\tgrok=c+a\log N+b\log(\eta\lambda)$ is undefined at
		$\eta\lambda=0$, so the no-decay column cannot enter the fit. With
		$\eta\lambda\in\{0,10^{-3},10^{-2},10^{-1}\}$ that leaves \emph{three}
		levels over two decades to estimate the coefficient carrying the competing
		hypothesis.
	\item Consequence if unfixed: the interval on $\hat b$ will likely admit both
		$0$ and $-1$, and the head claim fails to discriminate for want of design,
		not for want of an effect.
	\item Fix (i): refine the decay axis to
		$\eta\lambda\in\{10^{-3},3\!\times\!10^{-3},10^{-2},3\!\times\!10^{-2},10^{-1}\}$,
		paying for it by dropping one width --- $N$ already spans $1.5$ decades and
		is not the contested axis.
	\item Fix (ii): report $\eta\lambda=0$ separately rather than forcing it into
		the log fit with an arbitrary offset.
	\item That column is then the most informative single cell: a weight-decay
		clock $\tgrok\propto(1-\beta)/(\eta\lambda)$ \citep{kim2026clock} predicts
		$\tgrok\to\infty$ as $\eta\lambda\to0$, yet the Tier~0 baseline
		\citep[App.~8.3]{kumar2024grokking} groks at finite time with no weight
		decay at all.
	\item Make it quantitative with the reciprocal form
		\begin{equation}
			1/\tgrok = c_0 + c_1(\eta\lambda) + c_2\log N,
			\label{eq:recip}
		\end{equation}
		which admits $\eta\lambda=0$ as the intercept and so uses every cell. A pure
		weight-decay clock requires $c_0=0$; nonzero $c_0$ bounds how much of the
		timing weight decay can account for, independently of $\hat b$.
	\item Report \eqref{eq:recip} \emph{alongside} the log fit, not instead of it:
		they answer different questions and disagreement between them is
		diagnostic. \TODO{$\hat c_0$ with CI; state whether it excludes zero}
\end{itemize}

\subsection{Tier 0: reproduction}
\begin{itemize}
	\item \TODO{figure: our $\alpha$ sweep next to \citet{kumar2024grokking}
		App.~8.3}
	\item State agreement/disagreement quantitatively; if $p=23$ at $90\%$ does
		not grok cleanly (only 529 pairs), say so and report the fallback used.
\end{itemize}

\subsection{The confound, demonstrated (H0)}
\begin{itemize}
	\item \TODO{figure: two cells with near-equal $D=\Fnorm{K_t-K_0}/\Fnorm{K_0}$
		but opposite mechanism --- one where $A$ climbs, one where $A$ is flat while
		$S$ collapses under weight decay}
	\item This is the one contribution that cannot fail; it follows from
		\eqref{eq:decomp} and costs nothing extra to produce.
\end{itemize}

\subsection{Tier 1: is there a threshold $\Astar$? (H1, H2)}
\begin{itemize}
	\item H1 test: for every cell record $A$ at the moment of the test-loss fall;
		report the coefficient of variation of $\Astar$ across cells with bootstrap
		CIs over seeds. \TODO{numbers}
	\item Kill criterion: $\Astar$ varying by more than a factor of $\approx2$
		across cells means no universal threshold --- report as a negative result;
		it still says alignment is not a sufficient statistic for timing.
	\item H2 test: censored regression of $\log\tgrok$ on an alignment-crossing
		time and a scale summary; H2 predicts the alignment term absorbs the scale
		term. Report partial correlations with bootstrap CIs. Kill criterion: a
		scale term that stays significant with $A$ in the model.
	\item \TODO{figure: $S_t$, $R_t$, $A_t$ and the two losses on one log-$t$ axis
		for a representative cell}
\end{itemize}

\subsection{Tier 2: the $(a,b)$ fit (head claim, H3)}
\begin{itemize}
	\item Fit $\log\tgrok=c+a\log N+b\log(\eta\lambda)$ on the crossed grid with
		bootstrap intervals; include the interaction term.
	\item Predictions to compare against: shape clock $a>0,b\approx0$
		($a=1$ under $O(1/N)$); weight-decay clock $a\approx0,b\approx-1$
		\citep{kim2026clock}; both nonzero $\Rightarrow$ two clocks, slower one
		operative.
	\item Direction check: wider $\Rightarrow$ lazier $\Rightarrow$ slower to
		reshape $\Rightarrow$ longer $\tgrok$. A negative $\hat a$ signals a
		parameterisation error, not a result.
	\item Report residuals where the power law fails. \TODO{$\hat a,\hat b$, CIs,
		figure}
\end{itemize}

\subsection{Tier 3: clamped norm under MSE}
\begin{itemize}
	\item Replicate $\tgrok\propto e^{a\rho}$ \citep{truong2026norm} under MSE;
		compare the fitted rate to their $\approx7.5$ (MLP, CE).
	\item Test whether $\Astar$ is invariant to $\rho$. \TODO{numbers or
		``not run''}
\end{itemize}

\subsection{Generality and robustness}
\begin{itemize}
	\item $p\in\{23,29,31\}$ and train fractions $\{50,70,90\}\%$: does $\Astar$
		survive? \TODO{cut first if compute is tight --- buys generality, not
		discrimination}
	\item Robustness: uncentred alignment, first-logit eNTK, $\tau$ sensitivity,
		and the training and test parts of the probe set.
\end{itemize}

\section{Conclusion}

\begin{itemize}
	\item What the decomposition buys: \TODO{one paragraph, tied to whichever of
		H1--H3 survived}
	\item Limitations to state explicitly: one task (modular addition), one
		architecture family (two-layer MLP), gradient-flow theory applied to
		large-$\eta$ discrete GD, alignment is one of several plausible timing
		variables, $\tgrok$ depends on threshold choices (mitigated by the
		sensitivity analysis), censoring model assumptions.
	\item Reproducibility: synthetic data, one script, fixed seeds, config in one
		table, code zip submitted. Say this explicitly.
	\item Future work: \TODO{e.g. attention architectures, CE replication at full
		grid, Tier~4 theory if not completed}
\end{itemize}

\section*{Contribution}

\TODO{confirm against what was actually done before submission}


\section*{Acknowledgement}

\begin{itemize}
	\item \TODO{AI-use statement --- required somewhere by the brief; decide
		proposal footnote vs.\ here}
	\item \TODO{compute resources (Colab/Kaggle/Bunya) if applicable}
\end{itemize}

\section*{Appendix A.}
\label{sec:appendix}

\TODO{Decide what to put here}
% At most 5 pages. Key messages stay in the main text; everything here must be
% referenced from it.
\begin{itemize}
	\item A.1 Derivation of \eqref{eq:decomp} in full: expand
		$\Fnorm{K_t-K_0}^2=\Fnorm{K_t}^2-2\langle K_t,K_0\rangle+\Fnorm{K_0}^2$ and
		substitute $\langle K_t,K_0\rangle=\Fnorm{K_t}\Fnorm{K_0}(1-R_t)$.
	\item A.2 Full configuration table and the pre-registered threshold values.
	\item A.3 Censoring model and the estimator used, with diagnostics.
	\item A.4 Robustness: uncentred alignment, first-logit eNTK, $\tau$
		sensitivity grids.
	\item A.5 $y^\top K^{-1}y=\lVert w-w_0\rVert^2$ near the kernel limit, and its
		use as a linearisation-breakdown predictor.
	\item A.6 Tier~4 theory: bound on $\lVert T_t\rVert$ for the quadratic-activation
		net, analytic $A_\infty$, and the resulting $\tgrok\sim A_\infty/\dot A$
		estimate. \TODO{include only if completed}
	\item A.7 Additional figures not needed for the main claims.
\end{itemize}

\smaller
\printbibliography

\newpage
Indicate whether all your group members consent for your report to be used as a
teaching resource by including one of the statement below.

\begin{center}
\emph{We give consent for this to be used as a teaching resource.}\\
\emph{We DO NOT give consent for this to be used as a teaching resource.}
\end{center}

\end{document}
